\documentclass[11pt]{article}

\usepackage[margin=1in]{geometry}
\usepackage{amsmath}
\usepackage{amssymb}

\usepackage{graphicx}
\usepackage{bm}
\usepackage{booktabs}
\usepackage{array}
\usepackage{amsthm}
\usepackage{enumerate}
\usepackage[numbers,sort&compress]{natbib}
\usepackage{hyperref}
\usepackage{lineno}

\hypersetup{
  colorlinks=true,
  linkcolor=blue,
  citecolor=blue,
  urlcolor=blue
}

\newtheorem{theorem}{Theorem}
\newtheorem{proposition}{Proposition}

\begin{document}

\title{Topology, Chirality, and Vacuum Density: Conditional CDFT Constructions}

\author{Steve Bico Mujjabi, MD\thanks{msbico@gmail.com}\\
VuraLabs, Kampala, Uganda\\
Independent researcher\\
ORCID: \href{https://orcid.org/0009-0001-0556-5516}{0009-0001-0556-5516}\\
\url{https://github.com/bampita-bico/CDFD-Part-I-Release}}

\date{July 2026}

\maketitle

\begin{abstract}
This paper records conditional algebra and modelling hypotheses used by the
first CDFT papers. It does not close their physical open problems. The stiffness
relation is a rearrangement of Paper I's calibrated equilibrium equation, not an
independent derivation. The standard classification of $T(2,n)$ knots remains
usable, but the Faddeev--Niemi soliton-energy result has not been shown to apply
to the CDFT vortex functional. Consequently it does not select a trefoil or a
particle family here.

The displayed vacuum-density expression back-solves $\rho_0$ from a mass scale
already fitted to observed lepton masses. The chirality interaction is a
symmetry-allowed ansatz whose coefficient and phase are not derived from a
vacuum equation of state. These constructions are retained as hypotheses and
conditional calculations, not physical predictions.

Across the series, the public CDFD Runtime expresses this equilibrium vacuum
limit as a reusable flow-constraint state grammar. The calculations below use
that equilibrium limit only; adaptive departures are left for later dynamical
work rather than fitted in this manuscript. The paper-local supplementary
scripts remain the audited reproduction layer for the figures and tables in
this Part I archive.
\end{abstract}

\noindent\textbf{Keywords:} Constraint-Driven Field Theory, Constraint-Driven Flux Dynamics, adaptive operating ratio, vacuum structure, Mujjabi tests, reproducible physics, topology, chirality, vacuum density

\vspace{0.75em}
\linenumbers

\section*{Project Context}
This manuscript constitutes Part I (Paper 03) of the multi-part series
\textit{Constraint-Driven Flux Dynamics (CDFD)}. The underlying mathematical
architecture builds directly upon the concepts established in Paper 02 of this
series~\cite{MujjabiII2026}.

\section*{Symbol Conventions}

\noindent\textbf{CDFD physical-vacuum symbol convention.}
\begin{center}
{\small
\begin{tabular}{p{0.16\linewidth}p{0.33\linewidth}p{0.41\linewidth}}
\hline
Symbol & Part I meaning & Physical reading \\
\hline
$J$ or $\Phi$ & Driving flux or throughput & Transport current, field throughput, circulation drive, or generic flux; a subscript fixes the exact channel \\
$C$ & Active constraint or capacity burden & Vacuum transport capacity, geometric confinement, stiffness, boundary resistance, or load-bearing limit \\
$S$ & Surface responsiveness & Local ability of the vacuum or modeled surface to reconfigure under drive \\
$M_s$ & Structural memory & Retained geometry, phase history, stiffness history, or accumulated constraint state \\
$\Psi_s$ & Adaptive operating ratio $(J/C)\cdot S\cdot M_s$ or $(\Phi/C)\cdot S\cdot M_s$ & Local bookkeeping gauge for constrained operation; it is not a new universal constant \\
$R,a,\chi$ & Ring radius, core radius, and aspect ratio $R/a$ & Geometry used to connect vortex structure to dimensionless electromagnetic scales \\
$\rho_0,c,\kappa$ & Vacuum density scale, propagation speed, and transport stiffness & Medium parameters used in stability, pressure, and self-consistency calculations \\
\hline
\end{tabular}
}
\end{center}

\noindent\textbf{Greek-symbol convention.}
\begin{center}
{\small
\begin{tabular}{p{0.16\linewidth}p{0.33\linewidth}p{0.41\linewidth}}
\hline
Symbol & Local role & Interpretation rule \\
\hline
$\alpha$ & Fine-structure constant unless a local equation states otherwise & In the Part I physics papers, un-subscripted $\alpha$ means the electromagnetic fine-structure constant \\
$\beta,\gamma$ & Local response, stiffness, coupling, or correction coefficients & Equation-local coefficients; subscripts bind them to the named mode, channel, or approximation \\
$\theta,\phi,\Phi_c$ & Phase, orientation, or chirality angles & Used for torus-knot orientation, vacuum phase, or chiral phase geometry as defined in the nearby equation \\
$\delta,\epsilon$ & Perturbation, residual, tolerance, or small-offset scales & Local stability margins, numerical residuals, or experimental tolerances \\
$\lambda,\mu,\nu$ & Rate, mass, mode, or frequency labels & Meaning is fixed by the equation, subscript, and physics context where the symbol appears \\
$\rho,\sigma,\tau$ & Density, spread/noise, or time-scale terms & Local density, variance, stress, relaxation, or transport-memory quantities \\
$\Omega,\Lambda$ & State/domain space or integrated measure where defined & Reserved for explicitly defined state spaces, spectra, or synthesis-level quantities \\
\hline
\end{tabular}
}
\end{center}


\section{Introduction}

Paper~I records a calibrated toy energy functional; it did not derive
$\alpha$ or a vacuum stiffness. Paper~II records a selected $\mathbb{Z}_3$
parametrisation; it did not establish particle topology. Both leave physical
open problems. This paper distinguishes conditional algebra from the still
unsupported physical mappings.

We proceed problem by problem, in order of increasing depth:
Problem~1 ($\S$\ref{sec:kappa}), Problem~4 ($\S$\ref{sec:topology}),
Problem~3 ($\S$\ref{sec:M}), Problem~2 ($\S$\ref{sec:theta}).
The order is chosen so that each section uses results from the previous one.
Section~\ref{sec:summary} records the remaining open physical questions.

\subsection{Constraint history, chirality, and charge}

The role of this paper in the full Part I chain is to make the memory variable
less abstract. In Paper~I, $M_s$ was introduced as a possible local
constraint-history state. Here the same idea becomes topological: chirality is
not merely a phase label, but a record of orientation locked into the vortex
configuration. This gives the first concrete reading of the Mujjabi Vacuum
Memory Law.

The source notes also point to a charge interpretation. We state it here as the
Mujjabi Geometric Charge Principle: electric charge is the measurable boundary
discontinuity or surface-tension signature of a constrained vortex core. This is
not yet a completed derivation of $e$. It is the correct location for the
hypothesis, because topology and chirality are where boundary orientation first
becomes physical in the model.


\section{Conditional Algebra for $\kappa$}
\label{sec:kappa}

Paper~I introduced the transport stiffness $\kappa$ and noted that it was
computed from the condition $\chi = 1/\alpha$ rather than derived from the
vacuum medium properties independently.
It remains unresolved physically because the target aspect ratio was used to
fit $\kappa$ in Paper~I.

\begin{proposition}
  Conditional on the selected equilibrium equation and chosen $\chi$ and
  $\beta$, the fitted parameter $\kappa$ obeys:
  \begin{equation}
    \kappa \;=\; \chi^2\bigl(\ln(8\chi) - \beta + 1\bigr)\,.
    \label{eq:kappa_closed}
  \end{equation}
\end{proposition}

\begin{proof}
  The equilibrium condition $dE_\text{norm}/d\chi = 0$ (Paper~I, Eq.~(5)) gives
  \begin{equation}
    \ln(8\chi) - \beta + 1 \;=\; \frac{\kappa}{\chi^2}\,.
  \end{equation}
  Multiplying both sides by $\chi^2$ gives Eq.~\eqref{eq:kappa_closed}.
  This is a rearrangement of the assumed equilibrium condition. It does not
  determine $\kappa$ independently or derive $\chi$.
\end{proof}

Numerically:
\begin{align}
  \kappa &= (137.036)^2 \times \bigl(\ln(8 \times 137.036) - 1.75 + 1\bigr)
  \nonumber \\
         &= 18{,}778.87 \times 6.2497 \;=\; 117{,}362.0\,.
  \label{eq:kappa_numerical}
\end{align}
This returns the value used to calibrate Paper~I. Agreement to machine precision
is a property of the same algebra, not an independent verification.


\section{Topological Context and Its Modelling Boundary}
\label{sec:topology}

Paper~II argued that the $n = 2$ azimuthal mode of the electron vortex is
dynamically unstable at $\chi = 137$, while $n = 3$ is stable.
This argument was an energy-scaling estimate and not fully rigorous.
Topology supplies a classification, not a selection argument for this model.
The idea that particles may be understood as topological solitons has
a long history, from Skyrme's baryon model~\cite{Skyrme1961} to
Ra\~{n}ada's knotted electromagnetic fields~\cite{Ranada1989}
and the Faddeev--Niemi knot solitons~\cite{FaddeevNiemi1997}.
The CDFT vacuum itself is conceived as a physical medium in the sense
of Dirac~\cite{Dirac1951} and Volovik~\cite{Volovik2003}, where
topological defects carry quantised charges.

\subsection{Torus knots and topological locking}

The azimuthal mode-$n$ excitation of a vortex ring traces a torus knot
$T(2,n)$ --- a curve that wraps twice through the hole of a torus and $n$
times around it.
The classification of $T(2,n)$ structures is well-known in knot theory~\cite{Rolfsen1976}:

\begin{itemize}
  \item $n = 1$: the unknot. The unperturbed electron vortex (Paper~I).
  \item $n$ even: a \emph{two-component link}, not a knot. The two components
    can be separated without cutting. There is no topological barrier to
    unravelling.
  \item $n$ odd, $n \geq 3$: a true \emph{knot}. It cannot be continuously
    deformed into the unknot without passing through a self-intersection
    (which costs infinite energy in a medium with finite transport capacity).
    Such excitations are \emph{topologically locked}.
\end{itemize}

Table~\ref{tab:torus_knots} lists the first seven torus knots $T(2,n)$ with
their classifications.

\begin{table}[h]
\caption{Torus knots $T(2,n)$: standard classification and unvalidated CDFT associations.}
  \label{tab:torus_knots}
  \begin{tabular}{clllll}
    \toprule
    $n$ & $T(2,n)$ & Knot? & Lobes & Symmetry & Physical role \\
    \midrule
    1 & Unknot              & No  & 1 & $\mathrm{U}(1)$ & No particle assignment \\
    2 & Hopf link           & No  & 2 & $\mathbb{Z}_2$  & Standard link classification \\
    3 & Trefoil             & \textbf{Yes} & 3 & $\mathbb{Z}_3$ & Candidate CDFT association \\
    4 & Solomon's link      & No  & 4 & $\mathbb{Z}_4$  & Standard link classification \\
    5 & Cinquefoil          & Yes & 5 & $\mathbb{Z}_5$  & No particle assignment \\
    6 & 3-component link    & No  & 6 & $\mathbb{Z}_6$  & Standard link classification \\
    7 & $T(2,7)$ knot       & Yes & 7 & $\mathbb{Z}_7$  & No particle assignment \\
    \bottomrule
  \end{tabular}
\end{table}

\subsection{Why $n = 3$ and not $n = 5, 7, \ldots$}

The Faddeev--Niemi result~\cite{FaddeevNiemi1997} gives a bound for a particular
Skyrme-type field theory, in which the minimum energy of a knotted soliton with Hopf charge $Q$ satisfies
$E_\text{min}(Q) \sim Q^{3/4}$.
The CDFT vortex functional has not been shown to be that field theory, nor has
its torus-knot index been shown to equal the Hopf charge. The following ratio is
therefore only a hypothetical analogue, not a CDFT energy result:
\begin{equation}
  \frac{E_\text{min}(T(2,n))}{E_\text{min}(T(2,3))}
  \;=\; \left(\frac{n}{3}\right)^{3/4}\,.
  \label{eq:energy_ratio}
\end{equation}
Within the cited model the trefoil is its lowest-charge nontrivial example. This
does not select a CDFT trefoil or establish a $47\%$ CDFT energy difference.

Table~\ref{tab:energy_ordering} shows the energy ordering.
\begin{table}[h]
\caption{Illustrative Faddeev--Niemi scaling, not an energy ordering for CDFT vortices.}
  \label{tab:energy_ordering}
  \begin{tabular}{llcc}
    \toprule
    Knot & Hopf charge $Q$ & $E/E_\text{trefoil}$ & Knot? \\
    \midrule
    $T(2,3)$ trefoil    & 3 & 1.000 & Yes \\
    $T(2,5)$ cinquefoil & 5 & 1.467 & Yes \\
    $T(2,7)$            & 7 & 1.888 & Yes \\
    $T(2,9)$            & 9 & 2.280 & Yes \\
    \bottomrule
  \end{tabular}
\end{table}

\begin{figure}[ht]
  \centering
  \includegraphics[width=0.8\textwidth]{figures/fig3_0_panel.pdf}
  \caption{Faddeev-Niemi scaling showing the energy of knotted solitons relative to the trefoil ($n=3$). Higher torus knots cost significantly more energy.}
  \label{fig:faddeev_niemi}
\end{figure}

\subsection{Why the model gives three charged-lepton modes}

The trefoil $T(2,3)$ has exactly $\mathbb{Z}_3$ symmetry --- three lobes
separated by $120^\circ$.
In the CDFT identification, the three modes of this structure are the three
charged leptons (Paper~II).
There is no `fourth mode' of the trefoil; the $\mathbb{Z}_3$ group has
exactly three elements.

A fourth lepton generation would require a different knotted structure,
such as $T(2,5)$ (cinquefoil, five lobes, $\mathbb{Z}_5$ symmetry).
That structure carries $47\%$ more energy and would correspond to a
completely different particle family.
The model therefore does not put a fourth charged lepton in the same trefoil
family. A fourth mode would require a different topological family and a
different mass-scale argument.

\textbf{Problem~4 is addressed inside the model.}
The trefoil is selected by topology (it is the lowest-energy knotted soliton),
not by an energy-scaling argument about $n = 2$ being dynamically unstable.
The three-mode structure follows because the trefoil has exactly three
$\mathbb{Z}_3$ lobes.


\section{A Conditional Parametrisation of $M$}
\label{sec:M}

The Brannen parametrisation (Paper~II) writes the lepton masses as
\begin{equation}
  m_k \;=\; M \bigl(1 + \sqrt{2}\cos(\theta + \tfrac{2\pi k}{3})\bigr)^2\,,
  \quad k = 0, 1, 2\,.
  \label{eq:brannen_again}
\end{equation}
The scale $M$ sets the overall mass of the lepton family.
Paper~II fitted $M \approx 313.9$~MeV but did not derive it.

\subsection{Decomposition of $M$}

In CDFT the vortex ring has a normalised energy $E_\text{norm}$ that depends
only on $\chi$ and $\beta$:
\begin{equation}
  g(\chi, \beta)
  \;\equiv\; E_\text{norm}(\chi)
  \;=\; \chi\bigl(\ln(8\chi) - \beta\bigr) + \frac{\kappa}{\chi}\,.
  \label{eq:g_factor}
\end{equation}

\begin{figure}[ht]
  \centering
  \includegraphics[width=0.8\textwidth]{figures/fig3_2_panel.pdf}
  \caption{The normalised CDFT vortex energy $E_{\mathrm{norm}}$ near equilibrium. The geometry $\chi_{\mathrm{eq}}$ is governed purely by this local minimum.}
  \label{fig:e_norm_equilibrium}
\end{figure}

The \emph{dimensional} vortex energy is
\begin{equation}
  E_\text{vortex} \;=\; 2\pi \rho_0 c^2 a^3 \cdot g(\chi, \beta)\,,
  \label{eq:E_dimensional}
\end{equation}
where $\rho_0$ is a model parameter, $c$ is the speed of light, and $a$ is the
standard classical electron radius. This dimensional expression is an ansatz,
not a derived vortex energy.

The trefoil energy scale $M$ is proportional to $E_\text{vortex}$:
\begin{equation}
  M \;=\; \rho_0 c^2 a^3 \cdot g(\chi, \beta)\,.
  \label{eq:M_decomposed}
\end{equation}
(The factor $2\pi$ is absorbed into the definition of the Brannen scale.)

\subsection{Known factors}

From Papers~I--II:
\begin{align}
  \chi &= 1/\alpha = 137.035999177 \nonumber \\
  \beta &= 1.75 \nonumber \\
  \kappa &= 117{,}362.0 \quad \text{(from Section~\ref{sec:kappa})} \nonumber \\
  a &= 2.8179 \times 10^{-15}~\text{m} \nonumber \\
  g(\chi, \beta) &= 137.036 \times (6.9997 - 1.75) + 117{,}362.0/137.036
  = 1{,}575.83\,. \label{eq:g_numerical}
\end{align}

From the Brannen fit (Paper~II, using the smallest amplitude for the electron):
\begin{equation}
  M \;=\; \frac{m_e}{A_e^2}
  \;=\; \frac{0.51100~\text{MeV}}{(0.04035)^2}
  \;=\; 313.89~\text{MeV}\,,
  \label{eq:M_from_masses}
\end{equation}
where $A_e = 1 + \sqrt{2}\cos(\theta + 2\pi k_e/3)$ is the electron amplitude
factor evaluated at the smallest-amplitude mode.

\subsection{Solving for $\rho_0$}

Combining Eqs.~\eqref{eq:M_decomposed} and \eqref{eq:M_from_masses}:
\begin{equation}
  \rho_0
  \;=\; \frac{M}{c^2 a^3 g(\chi,\beta)}
  \;=\; \frac{313.89~\text{MeV}}{c^2 \cdot (2.818 \times 10^{-15})^3
        \cdot 1{,}575.83}
  \;=\; 1.587 \times 10^{13}~\text{kg/m}^3\,.
  \label{eq:rho0_numerical}
\end{equation}
This is a back-solved numerical value because $M$ was fitted from observed
lepton masses. Its numerical scale does not establish a physical vacuum density.


\section{A Symmetry-Allowed Chirality Ansatz}
\label{sec:theta}

The phase $\theta \approx 12.73^\circ$ was fitted in Paper~II but not derived.
The following is a symmetry-allowed ansatz; it does not derive a mechanism or
reduce the physical problem to knowledge of $\rho_0$.

\subsection{Vacuum chirality and $\mathbb{Z}_3$ breaking}

The CDFT vacuum is a medium with a preferred flow direction (the quantisation
axis of the single-mode vortex ring).
This background flow has a \emph{helicity} --- a correlation between the
velocity and vorticity fields --- which we parameterise by a phase $\phi_c$,
the \emph{vacuum chirality phase}.

The trefoil vortex, placed in this background, acquires an interaction energy:
\begin{equation}
  E_\text{chiral}(\theta)
  \;=\; E_{\mathbb{Z}_3} + \lambda_c \cos(3\theta + \phi_c)\,,
  \label{eq:E_chiral}
\end{equation}
where $E_{\mathbb{Z}_3}$ is the $\mathbb{Z}_3$-symmetric part (independent
of $\theta$, hence independent of the mass hierarchy) and $\lambda_c$ is the
coupling strength.

The $\cos(3\theta + \phi_c)$ form is required by the $\mathbb{Z}_3$ symmetry
of the trefoil: any interaction that respects $\mathbb{Z}_3$ must be periodic
in $\theta$ with period $2\pi/3$, and the lowest-order such interaction is the
fundamental $\cos(3\theta + \phi_c)$.

\subsection{Equilibrium condition for $\theta$}

\begin{proposition}
  Conditional on the assumed chiral energy, a stationary orientation $\theta$ satisfies
  \begin{equation}
    3\theta + \phi_c \;=\; n\pi\,,
    \quad n \in \mathbb{Z}\,.
    \label{eq:theta_equation}
  \end{equation}
\end{proposition}

\begin{figure}[ht]
  \centering
  \includegraphics[width=0.8\textwidth]{figures/fig3_1_panel.pdf}
  \caption{The chiral energy surrogate $\cos(3\theta + \phi_c)$ plotted against $\theta$. The fitted $\theta$ exactly minimises this interaction.}
  \label{fig:chiral_energy}
\end{figure}

\begin{proof}
  The equilibrium condition $dE_\text{chiral}/d\theta = 0$ gives
  $-3\lambda_c \sin(3\theta + \phi_c) = 0$,
  hence $3\theta + \phi_c = n\pi$.
  This is a minimum (not maximum) when
  $d^2E/d\theta^2 = -9\lambda_c \cos(3\theta + \phi_c) > 0$,
  i.e., when $n$ is odd and $\lambda_c > 0$.
  This does not establish a physical domain, unique physical solution, or the
  existence of the chiral interaction.
\end{proof}

\subsection{Computing $\phi_c$ from observation}

From Paper~II, the fitted phase is $\theta_\text{obs} = 12.7325^\circ$.
Equation~\eqref{eq:theta_equation} with $n = 1$ gives:
\begin{equation}
  \phi_c \;=\; \pi - 3\theta_\text{obs}
  \;=\; 180^\circ - 38.20^\circ \;=\; 141.80^\circ\,.
  \label{eq:phi_c}
\end{equation}
This is verified numerically in supplementary Table~4:
$\cos(3\theta_\text{obs} + \phi_c) = \cos(\pi) = -1$ to machine precision.

\subsection{Connection to $\rho_0$}

The vacuum chirality phase $\phi_c$ arises from the helicity of the vacuum
flow field.
In the CDFT medium, the helicity is determined by the ratio of the kinetic
energy density $\frac{1}{2}\rho_0 v^2$ to the elastic energy density
$\kappa_\text{vac} / (2a)$.
This ratio is a function of $\rho_0$, $c$, $a$, and the vacuum
compressibility --- all quantities that the vacuum equation of state relates
to $\rho_0$.

No derivation in this paper establishes $\phi_c=\phi_c(\rho_0)$. The displayed
phase is back-solved from a fitted $\theta$, and the stationary relation is only
conditional calculus for the assumed interaction.


\section{Scope Boundary and Conclusion}
\label{sec:summary}

This paper retains four conditional calculations, none of which closes a
physical problem. First, the $\kappa$ expression is a rearrangement of the
Paper I calibration equation. Second, the knot classifications are standard
topology, but their physical application and energy ordering have not been
derived for CDFT. Third, the displayed $\rho_0$ is back-solved from an observed
mass fit, so it is not a vacuum-density determination. Fourth, the stationary
condition for the assumed chiral cosine is calculus conditional on an ansatz;
the ansatz, phase, and coupling are not independently derived.

The open requirements are an independently specified physical energy functional,
a dimensional closure for its parameters, a justified relation to any soliton
model, and a prespecified discriminating observation. Until these exist, this
paper must not be cited as a derivation of $\alpha$, lepton masses, vacuum density,
trefoil selection, or a universal knot hierarchy.


\section*{Supplementary Material}

This manuscript is accompanied by release-archive supplementary material in
the canonical Part I tree. The relevant paper-local Python scripts, numerical checks, generated figure outputs,
tabulated data, figure assets, and environment notes are maintained under
\texttt{notebooks/}, \texttt{outputs/}, \texttt{figures/}, and
\texttt{REPRODUCIBILITY.md} in \texttt{Part\_I\_Fundamental\_Physics/}.


\section*{Data and Code Availability}

The release-archive supplementary material is included in the archived Part I
release on Zenodo at \url{https://doi.org/10.5281/zenodo.20250820}. The current version DOI is assigned by Zenodo after each GitHub release. Zenodo is the
permanent release archive, and the public source archive is available on GitHub at
\url{https://github.com/bampita-bico/CDFD-Part-I-Release}. The broader public CDFD
Runtime is separate reusable software infrastructure; the numerical values reported
here are reproduced from this paper's Supplementary Material.


\section*{Acknowledgments}
The author used AI-assisted tools for formatting checks and code-verification
support. The scientific claims, interpretations, and final manuscript decisions
are the author's responsibility.

The author is indebted to the foundational work of
L.~Faddeev and A.J.~Niemi~\cite{FaddeevNiemi1997} (knots and particles),
T.H.R.~Skyrme~\cite{Skyrme1961} (topological solitons),
A.F.~Ra\~{n}ada~\cite{Ranada1989} (knotted field solutions),
G.E.~Volovik~\cite{Volovik2003} (vacuum as condensate),
and D.~Rolfsen~\cite{Rolfsen1976} (torus knot classification)
--- whose work provides the mathematical and physical foundations
for the topological arguments in this paper.

\section*{Declarations}

\subsection*{Funding}
This research received no external institutional funding.

\subsection*{Author Contributions}
Steve Bico Mujjabi: conceptualization, methodology, formal analysis, software,
validation, visualization, writing--original draft, and writing--review and
editing.

\subsection*{Conflicts of Interest}
The author declares no conflicts of interest.

\section*{References}
\bibliographystyle{plainnat}
\bibliography{references}

\end{document}
